\documentclass[letterpaper,10pt,conference]{ieeeconf}
\overrideIEEEmargins
\usepackage[T1]{fontenc}
\usepackage{amsmath,amsfonts}
\usepackage{algpseudocode}
\usepackage{algorithm}
\usepackage{array}
\usepackage[caption=false]{subfig}
\usepackage{tabularx}
\usepackage{colortbl}
\usepackage{textcomp}
\usepackage{stfloats}
\usepackage{multirow}
\usepackage{verbatim}
\usepackage{graphicx}
\usepackage{subcaption}
\usepackage{cite}
\usepackage{flushend}
\usepackage{diagbox}
\usepackage[dvipsnames]{xcolor}
\usepackage{threeparttable}
\usepackage{booktabs}
\usepackage{empheq}
\usepackage{makecell}
\usepackage{url}
\usepackage{soul}
\usepackage{tikz}
\let\labelindent\relax
\usepackage{enumitem}
\usepackage{gensymb}
\usepackage{amssymb}
\usepackage[colorlinks=true, linkcolor=black, citecolor=black, urlcolor=blue]{hyperref}
\urlstyle{same} % makes the URL use the normal text font (not typewriter)
\definecolor{redvw}{HTML}{F36B9C}
\definecolor{bluevw}{HTML}{1E88E5}
\definecolor{greenvw}{HTML}{005648}
\definecolor{navyblue}{RGB}{0, 83, 137}

\newcommand{\victor}[1]{{\color{navyblue}#1}}
\newcommand{\victornote}[1]{{\color{orange}#1}}
\newcommand{\rp}[1]{{\color{blue}#1}}

\newcommand{\Estar}{E_{\star}}
\newcommand{\dd}{\mathrm{d}}
\newcommand{\ind}{\mathbf{1}}



\title{\bfseries Hierarchical Energy-Work Control for Rotary-Pendulum Capture}
\author{}




\begin{document}
\maketitle
\thispagestyle{empty}
\pagestyle{empty}

\section{Introduction}
% Research objective and the role of intervention timing and prediction horizon.
Spacecraft and marine navigation motivate path-planning and control strategies that use natural dynamical flows to passively coast an agent toward its goal, maneuvered by occasional control interventions. The research goal is a act--coast controller that completes a navigation task with a limited number of short actuation periods separated by longer un-actuated motion. Each intervention must shape the subsequent coasting trajectory, so its timing matters alongside the applied control~\cite{heemels2012}. The project aim to jointly select the control sequence, the actuation and coasting durations, and a state-dependent prediction horizon. This horizon determines how far ahead the controller predicts motion to assess an intervention, driven by limited computational cost.

% Energy-based nominal actions and the need for phase-dependent torque selection.
We use a simulated rotary pendulum as a benchmark, where a motor-driven horizontal arm must swing an free-hinged pendulum at the arm's tip to upright. Building on energy-based swing-up~\cite{astrom2000}, we derive a heuristic from the mechanical dynamics to propose nominal control actions. It represents the physical state by a reduced order energy state: the arm kinetic energy, pendulum kinetic energy, and gravitational potential energy. Their sum determines the energy deficit or surplus relative to upright rest and hence the requested motor work. Matching the energy state alone does not lead to correct motion, since multiple physical states can have the same energy state. The controller therefore uses the phase information to derive the correct torque to deliver requested work. This heuristic provides a starting prior for act--coast development, with its proposed nominal actions guide the later MPC prediction and reinforcement learning training.

\section{Problem Formulation}
% Geometry, angular conventions, and indirect pendulum actuation.
A rotary pendulum system contains a motor-driven horizontal arm carring an unactuated pendulum hinged at its tip, which swings in the vertical plane perpendicular to the arm. The unwrapped arm angle $\theta\in\mathbb R$ measures rotation about the vertical axis from a fixed horizontal reference, retaining complete revolutions. The pendulum angle $\alpha\in[-\pi,\pi)$ is zeroed from downward, with upright represented by $\pm\pi$. Their signed angular velocities are $\omega=\dot\theta$ and $\nu=\dot\alpha$, with derivatives taken before angular wrapping. Motor torque $u$ is positive toward increasing $\theta$.

% Observed state, admissible input, initial conditions, and task objective.
At each control update $t_k$, for $k\in\mathbb{N}$, the controller observes the physical state $x$ and holds torque $u_k$ for state transition $x_k\rightarrow x_{k+1}$, with
\begin{equation}
\begin{aligned}
 x&=(\theta,\alpha,\omega,\nu)^{\mathsf T},\qquad |u(t)|\le U,\\
 x(0)&\in\mathcal X_0\subset\mathbb R\times[-\pi,\pi)\times\mathbb R^2,
\end{aligned}
\label{eq:state}
\end{equation}
where $t$ denotes time, $U>0$ bounds torque, and the bounded initial set $\mathcal X_0$ is sampled according to Table~\ref{tab:initial}. Define the upright error and target as
\begin{equation}
\begin{aligned}
 \beta&=\operatorname{wrap}(\alpha-\pi),\\
 \mathcal X_\star&=\{x:\beta=0,\ \omega=\nu=0,\ |\theta|<\theta_{\max}\}.
\end{aligned}
\label{eq:target}
\end{equation}
Here $\operatorname{wrap}$ selects the equivalent angle in $[-\pi,\pi)$, and $\theta_{\max}>0$ prescribes the tracking envelope. The task is swing the pendulum up to an upright rest, with constraints of torque bound and arm angle range.

% Acquisition, continued motion, and the cost of control interventions.
The upright goal requires $|\beta|\le\varepsilon_\alpha$ at consecutive time steps with positive tolerance $\varepsilon_\alpha$. The controller's performance is evaluated by success possibility within limited time. The trajectories are also evaluated within the time limit $\mathcal I=[t_0,t_0+T], T>0$ to compute
\begin{equation}
\begin{aligned}
 \beta_{\mathrm{rms}}&=\left(\frac{1}{T}\int_{\mathcal I}\beta(t)^2\,\dd t\right)^{1/2},\\
 \rho_\alpha&=\frac{1}{T}\int_{\mathcal I}\ind_{\{|\beta(t)|\leq\varepsilon_\alpha\}}\,\dd t.
\end{aligned}
\label{eq:holdingmetrics}
\end{equation}

Here $\ind$ equals one when its condition holds and zero otherwise. These two metrics measure root-mean-square (RMS) angular deviation of the pendulum from the goal, and the arm's deviation from neutral position while holding the pendulum upright. They assess the controller's ability to converge fast, hold upright steadily, and stay within operation envelope. We also define
\begin{equation}
\begin{aligned}
 W_{\mathrm{abs}}&=\int_{\mathcal I}|u(t)\omega(t)|\,\dd t,\\
 T_{\mathrm{act}}&=\int_{\mathcal I}\ind_{\{u(t)\neq0\}}\,\dd t.
\end{aligned}
\label{eq:resourcemetrics}
\end{equation}
where absolute shaft work $W_{\mathrm{abs}}$ integrates energy injected and removed, and $T_{\mathrm{act}}$ measures nonzero-actuation duration. Total computation time is also recorded to reflect planning cost. \victor{And later act--coast evaluation additionally counts times of interventions and measures zero-torque coast durations.} These metrics collectively form a systematic way to assess the performance of controllers on this rotary pendulum system.










\section{Method}
% Work selection and torque realization use complementary state information.
With bounded torque and arm travel, the rotary pendulum must raise an unactuated pendulum through configuration-dependent coupled swing-ups. The swing-ups require successive oscillations, so pursuing monotonically decreasing angular error is not suitable. We therefore propose a method to compute a desired energy change, then use the physical state to determine the energy distribution and the corresponding torque realization. Section~\ref{sec:dynamics} derives this energy representation; Section~\ref{sec:contorller} gives the work-selection and torque-decoding scheme.

\subsection{Rotary-pendulum dynamics and energy representation}
\label{sec:dynamics}
% The mechanical energy determines the coupled equations of motion.
Both arm and pendulum are modeled as uniform rigid rods with ideal joints, zero damping, and negligible pendulum axial inertia. Let $m_a,r$ denote arm mass and length, $m_p,l$ pendulum mass and length, and $g$ gravitational acceleration, all positive. The arm inertia about the motor axis is $I_a=m_ar^2/3$. The pendulum center-of-mass distance from the hinge is $c=l/2$, and the transverse inertia about its center of mass is $I_c=m_pl^2/12$. The parallel-axis theorem gives the pendulum's hinge inertia $J=I_c+m_pc^2$. Define the coupling coefficient $b=m_prc$ and gravitational energy scale $G=m_pgc$. For generalized coordinates $q=(\theta,\alpha)^{\mathsf T}$, the kinetic energy $K$, potential energy $V$, and Lagrangian $\mathcal L$ are
\begin{equation}
\begin{aligned}
 A(\alpha)&=I_a+m_pr^2+J\sin^2\alpha,\\
 C(\alpha)&=b\cos\alpha,\\
 K(q,\dot q)&=\tfrac12A\omega^2+C\omega\nu+\tfrac12J\nu^2,\\
 V(\alpha)&=G(1-\cos\alpha),\qquad \mathcal L=K-V.
\end{aligned}
\label{eq:lagrangian}
\end{equation}
Here $A$ is the inertia associated with arm rotation, including the pendulum contribution $J\sin^2\alpha$ about the motor axis. The term $C\omega\nu$ couples both motions, while $\frac{1}{2}J\nu^2$ describes pendulum rotation about its hinge. Potential energy is zero at downward rest. Applying the forced Euler--Lagrange relation $\frac{\dd}{\dd t}(\partial\mathcal L/\partial\dot q)-\partial\mathcal L/\partial q=(u,0)^{\mathsf T}$ and solving for both accelerations gives the following governing equation of the dynamics. Let primes denote differentiation with respect to $\alpha$, giving $A'=2J\sin\alpha\cos\alpha$, $C'=-b\sin\alpha$, and $V'=G\sin\alpha$, we have 
\begin{equation}
\begin{aligned}
 M(\alpha)&=\begin{bmatrix}A&C\\C&J\end{bmatrix},\\
 \dot x=f(x,u)&=
 \begin{bmatrix}
 \omega\\
 \nu\\
 M(\alpha)^{-1}\!
 \begin{bmatrix}
 u-A'\omega\nu-C'\nu^2\\
 \tfrac12A'\omega^2-V'
 \end{bmatrix}
 \end{bmatrix}.
\end{aligned}
\label{eq:dynamics}
\end{equation}
The zero pendulum component of the generalized input $(u,0)^{\mathsf T}$ expresses the absence of direct pendulum actuation. The positive-definite inertia matrix $M$ couples the two acceleration equations; its inverse determines $(\dot\omega,\dot\nu)^{\mathsf T}$, the last two components of $f$. The remaining terms account for velocity-dependent inertial torques and gravity. Integrating $f$ predicts the physical actuated trajectory of the rotary pendulum.

% The physical energy partition separates motion from elevation.
The physical state describes the oscillatory motion, but angular error alone does not provide suitable input to reach upright. Therefore, we use a reduced-order energy state $e$ to separate arm motion, pendulum motion, and elevation, and use its total-energy error to request motor work. Define the energy state $e$ and its goal state $e_\star$
\begin{equation}
\begin{aligned}
 e=H(x)&=(K_a,K_p,V)^{\mathsf T},\\
 e_\star=(0,0&,\Estar)^{\mathsf T},\\
 \text{s.t.} \quad K_a&=\tfrac12I_a\omega^2,\qquad K_p=K-K_a,\\
 E&=K_a+K_p+V,\\
 \Estar&=2G.
\end{aligned}
\label{eq:encoder}
\end{equation}
Here $H(x)$ maps the physical state $x$ to the energy state $e$. Arm kinetic energy $K_a$ excludes the pendulum; pendulum kinetic energy $K_p$ includes motion induced by the moving hinge and the coupling term. At the upright-rest target $e_\star$, total energy $E$ consists entirely of potential energy. Differentiating these energies along $x$ gives
\begin{equation}
\begin{aligned}
 \dot K_a&=I_a\omega\dot\omega,\\
 \dot K_p&=u\omega-I_a\omega\dot\omega-G\sin\alpha\,\nu,\\
 \dot V&=G\sin\alpha\,\nu,\qquad \dot E=u\omega.
\end{aligned}
\label{eq:powers}
\end{equation}
Motor power $u\omega$ changes total energy, while the other terms describe its exchange among the three components. Regulating $E$ alone leaves its distribution undetermined, so an ideal approach is to drive both $E$ and $V$ to $\Estar$, the remaining kinetic energy $K_a+K_p=E-V$ then vanishes to zero.

% The energy-map Jacobian relates physical motion to energy rates.
Express the energy dynamics governing equation through the chain rule,
\begingroup
\setlength{\arraycolsep}{4pt}
\begin{equation}
\begin{aligned}
 \dot e&=DH(x)f(x,u),\\
 DH(x)&=
 \begin{bmatrix}
 0&0&I_a\omega&0\\[3pt]
 0&\begin{gathered}\tfrac12A'\omega^2\\{}+C'\omega\nu\end{gathered}
 &\begin{gathered}(A-I_a)\omega\\{}+C\nu\end{gathered}&C\omega+J\nu\\[3pt]
 0&V'&0&0
 \end{bmatrix}.
\end{aligned}
\label{eq:energy-jacobian}
\end{equation}
\endgroup
The Jacobian $DH(x)=\partial H/\partial x$ converts physical-state rates into the energy rates in Eq.~\eqref{eq:powers}. Its rows correspond to $(K_a,K_p,V)$ and columns to $(\theta,\alpha,\omega,\nu)$. The zero first column reflects independence from absolute arm position, which remains necessary for constraint evaluation. If $H$ had a unique inverse, we could recover $x$ from $e$ and predict energy transfer using only $(e,u)$. However, with both velocities fixed, the states at $\alpha$ and $-\alpha$ have the same $e$ but opposite $\dot V$ when $\sin\alpha\,\nu\ne0$. Thus $e$ does not uniquely determine $x$, and $H$ has no unique inverse. The differing energy rates show that $(e,u)$ alone cannot determine $\dot e$. Determining the required torque for control therefore requires phase information to distinguish the motions. 













\subsection{Work selection and phase-dependent torque decoding}
\label{sec:contorller}
% Energy feedback requests work over a prescribed interval of bounded actuation.
The energy representation $H(x)$ in Section~\ref{sec:dynamics} specifies the required energy correction, while $\dot x=f(x,u)$ determines the response to bounded torque. At each update $t_k$, the controller observes $x_k=x(t_k)$ and holds $u(t)=u_k$ until $t_{k+1}$, with prescribed duration $\Delta=t_{k+1}-t_k>0$. The current energies $e_k=(K_{a,k},K_{p,k},V_k)^{\mathsf T}=H(x_k)$ give the proportional work request
\begin{equation}
\begin{aligned}
 E_k&=K_{a,k}+K_{p,k}+V_k,\\
 w_k&=\operatorname{clip}\!\left(\kappa(\Estar-E_k),-W_{\max},W_{\max}\right).
\end{aligned}
\label{eq:workpolicy}
\end{equation}
Here $E_k$ is current total energy, and $w_k$ is the signed work requested before the next update. The dimensionless gain $\kappa>0$ and request bound $W_{\max}>0$ are design choices. The operator $\operatorname{clip}$ restricts $\kappa(\Estar-E_k)$ to $[-W_{\max},W_{\max}]$. The sign of $w_k$ requests energy addition for $E_k<\Estar$ and removal for $E_k>\Estar$. The torque bound $|u|\le U$ alone does not ensure delivery of $w_k$ over $\Delta$.

% A requested total energy specifies a plane but leaves the energy distribution free.
Delivering $w_k$ fixes the next-step total energy at $E_k+w_k$ but leaves the energy distribution undetermined. For a candidate energy state $z=(z_a,z_p,z_V)^{\mathsf T}$, the equality $z_a+z_p+z_V=E_k+w_k$ defines a plane in three-dimensional energy space. The components $z_a,z_p,z_V$ denote arm kinetic, pendulum kinetic, and potential energy, respectively. Let $\mathcal E=\{H(x):x\text{ is a physical state}\}$ be the continuous region of physically realizable energies with unrestricted velocities. The mechanical bounds on $\mathcal E$ are given in Eqs.~\eqref{eq:energy-feasibility}--\eqref{eq:energy-inertias}. The physically realizable part of $z_a+z_p+z_V=E_k+w_k$ is
\begin{equation}
 \mathcal S(e_k,w_k)=\{z\in\mathcal E:z_a+z_p+z_V=E_k+w_k\}.
 \label{eq:surface}
\end{equation}
For $E_k+w_k>0$, $\mathcal S(e_k,w_k)$ is a continuous planar region. At $E_k+w_k=0$, $\mathcal S(e_k,w_k)$ reduces to the origin; for $E_k+w_k<0$, $\mathcal S(e_k,w_k)$ is empty. Increasing $z_V$ on $\mathcal S(e_k,w_k)$ reduces $z_a+z_p$. Membership in $\mathcal S(e_k,w_k)$ does not establish attainability from $x_k$ under $|u|\le U$.

% One continuous physical prediction determines the next-step energy for each torque.
Predicting attainability requires $x_k$ because $H(x_k)$ does not uniquely determine the signed motion, as shown in Section~\ref{sec:dynamics}. For a trial torque $u\in[-U,U]$, let $x_u(s)$ denote the predicted physical state after elapsed time $s\in[0,\Delta]$. Holding $u$ constant, $f(x,u)$ and $H(x)$ define
\begin{equation}
\begin{aligned}
 \frac{\dd x_u}{\dd s}(s)&=f(x_u(s),u),\qquad x_u(0)=x_k,\\
 H_\Delta(x_k,u)&=H\!\left(x_u(\Delta)\right).
\end{aligned}
\label{eq:energy-prediction}
\end{equation}
The continuous predictor $x_u(s)$ retains configuration and signed velocities, while $H_\Delta(x_k,u)$ expresses the predicted next-step state in energy coordinates. The Jacobian $DH(x)=\partial H/\partial x$ in Eq.~\eqref{eq:energy-jacobian} gives the instantaneous rates $\dot e=DH(x)f(x,u)$. Equation~\eqref{eq:energy-jacobian} therefore justifies computing $H_\Delta(x_k,u)$ from $x_u(\Delta)$ without separately integrating $\dot e$.

% The work balance connects predicted arm displacement to the requested energy change.
To determine whether $u$ delivers $w_k$, compute the signed work $W_\Delta(x_k,u)$ from $x_u(s)$. Let $\theta_u(s)$ and $\omega_u(s)$ be the unwrapped arm angle and angular velocity in $x_u(s)$. Since $E=K_a+K_p+V$, Eq.~\eqref{eq:powers} gives $\dot E=\dot K_a+\dot K_p+\dot V=u\omega$ at each continuous time. Integrating $u\omega_u(s)$ with respect to elapsed time $s$ yields
\begin{equation}
\begin{aligned}
 W_\Delta(x_k,u)&=\int_0^\Delta u\,\omega_u(s)\,\dd s\\
 &=u\bigl[\theta_u(\Delta)-\theta_k\bigr],\\
 E\!\left(x_u(\Delta)\right)&=E_k+W_\Delta(x_k,u).
\end{aligned}
\label{eq:work}
\end{equation}
Here $E(x)=K_a(x)+K_p(x)+V(x)$ is total energy at $x$, and $\theta_k$ is the observed arm angle. In Eq.~\eqref{eq:work}, $u$ is constant and $\omega_u(s)=\dd\theta_u(s)/\dd s$, giving the displacement $\theta_u(\Delta)-\theta_k$. Changing $u$ also changes $\theta_u(\Delta)$ through Eq.~\eqref{eq:dynamics}, so $W_\Delta(x_k,u)$ depends on coupled motion. At rest, $u\omega_u(0)=0$, although $u$ can initiate motion. Predicting $W_\Delta(x_k,u)$ includes subsequent energy transfer without dividing by $\omega_u(0)$.

% Varying the scalar torque produces a continuous family of attainable energies.
Each $u\in[-U,U]$ now has a predicted next-step energy $H_\Delta(x_k,u)$. For fixed $x_k$ and $\Delta$, varying $u$ produces
\begin{equation}
 \mathcal C_{x_k}=\{H_\Delta(x_k,u):u\in[-U,U]\}.
 \label{eq:curve}
\end{equation}
Smooth dependence of $f(x,u)$ and $H(x)$ on state and input makes $\mathcal C_{x_k}$ a continuous parameterized curve in energy space. Each $H_\Delta(x_k,u)$ belongs to $\mathcal E$ because $H_\Delta(x_k,u)=H(x_u(\Delta))$. Equation~\eqref{eq:work} gives $E(x_u(\Delta))=E_k+W_\Delta(x_k,u)$. Hence $H_\Delta(x_k,u)\in\mathcal S(e_k,w_k)$ if and only if
\[
 W_\Delta(x_k,u)=w_k.
\]
Solving $W_\Delta(x_k,u)=w_k$ identifies $\mathcal C_{x_k}\cap\mathcal S(e_k,w_k)$ and recovers $u$ directly as the parameter of $\mathcal C_{x_k}$, without inverting $H$. Multiple solutions $u$ can yield different $H_\Delta(x_k,u)$ and arm excursions. If $W_\Delta(x_k,u)=w_k$ has no solution in $[-U,U]$, no admissible constant torque delivers $w_k$ over $\Delta$.

% Numerical integration and scalar root refinement supply the candidate torques.
A numerical solution of $W_\Delta(x_k,u)=w_k$ requires approximating $x_u(s)$ in Eq.~\eqref{eq:energy-prediction}. Divide $\Delta$ into $N$ equal integration steps of length $h=\Delta/N$, with prescribed integer $N\ge1$. For each $u$, let $\widehat x_j$ approximate $x_u(jh)$, where $j=0,\ldots,N$ indexes integration samples. The angular components of $\widehat x_j$ are $\widehat\theta_j$ and $\widehat\alpha_j$. Using fourth-order Runge--Kutta integration, denoted $\operatorname{RK4}$, the discrete prediction is
\begin{equation}
\begin{aligned}
 \widehat x_0&=x_k,\\
 \widehat x_{j+1}&=\operatorname{RK4}(\widehat x_j,u;h),\qquad j<N,\\
 \widehat x_N&\approx x_u(\Delta),\\
 H(\widehat x_N)&\approx H_\Delta(x_k,u),\\
 \widehat W_\Delta(x_k,u)&=\sum_{j=0}^{N-1}u(\widehat\theta_{j+1}-\widehat\theta_j)\\
 &=u(\widehat\theta_N-\theta_k).
\end{aligned}
\label{eq:state-prediction}
\end{equation}
Hats distinguish numerical predictions from the continuous solution $x_u(s)$. After $N$ integration steps, $H(\widehat x_N)$ approximates $H_\Delta(x_k,u)$ at the next update. The discrete work $\widehat W_\Delta(x_k,u)$ adds work over angular increments; constant $u$ makes the finite sum telescope to $u(\widehat\theta_N-\theta_k)$. Numerical integration introduces error into the energy balance, so $E(\widehat x_N)\approx E_k+\widehat W_\Delta(x_k,u)$.

% Bracketing and bisection approximate intersections without an analytic torque inverse.
To approximate $W_\Delta(x_k,u)=w_k$, define the numerical work residual $R(u)$ and a uniform grid of $n_u$ torques $u_i$, with integer $n_u\ge2$,
\begin{equation}
\begin{aligned}
 R(u)&=\widehat W_\Delta(x_k,u)-w_k,\\
 u_i&=-U+\frac{2Ui}{n_u-1},\qquad i=0,\ldots,n_u-1.
\end{aligned}
\label{eq:torque-search}
\end{equation}
The equation $R(u)=0$ is solved numerically by bisection, without a closed-form inverse for $u$. For $i=0,\ldots,n_u-2$, continuity of $R(u)$ and $R(u_i)R(u_{i+1})\le0$ ensure a root in $[u_i,u_{i+1}]$. At each iteration, compute $R(u)$ at the bracket midpoint and retain a half whose boundary values of $R(u)$ have nonpositive product. After $L$ iterations, with prescribed integer $L\ge0$, the bracket width is $2U/[(n_u-1)2^L]$. Retain the bracket midpoints and grid torques in $\mathcal U_{\mathrm{test}}\subseteq[-U,U]$, excluding nonfinite predictions. Predictions $H(\widehat x_N)$ for $u\in\mathcal U_{\mathrm{test}}$ form a finite collection of points approximating $\mathcal C_{x_k}$. Roots without a sign change in $R(u)$, or multiple roots between adjacent grid torques, can be missed.

% The integration and search counts determine the arithmetic cost per control update.
For the fixed four-state dynamics $f(x,u)$, computing $R(u)$ requires $N$ Runge--Kutta steps, or $4N$ evaluations of $f(x,u)$. The implementation batches all $n_u-1$ grid intervals during each of $L$ bisection iterations. Including the initial grid and prediction of at most $2n_u-1$ candidate torques gives $\mathcal O(n_uN(L+1))$ arithmetic work per update. Batching evaluates trial torques in parallel; elapsed computation time also depends on hardware.

% Candidate ranking balances energy distribution, work agreement, and arm excursion.
The condition $R(u)=0$ expresses work agreement but does not rank $u\in\mathcal U_{\mathrm{test}}$ by $H(\widehat x_N)$ or $\widehat\theta_j$. Increasing $V$ toward $\Estar$ complements the feedback $w_k$ because $K_a+K_p=E-V$. Since $H(x)$ omits $\theta$, arm excursion requires checking $\widehat\theta_j$ for all $j=0,\ldots,N$. In particular, $|\widehat\theta_N|\le\theta_{\max}$ does not exclude $|\widehat\theta_j|>\theta_{\max}$ at an earlier sample. Let $\widehat V^+(x_k,u)$ denote predicted next-step potential energy and $\widehat\Theta(x_k,u)$ the largest sampled absolute arm angle. The score $\mathcal J(x_k,u;w_k)$ combines $\widehat V^+(x_k,u)$, $R(u)$, and $\widehat\Theta(x_k,u)$ to select $u_k$ through
\begin{equation}
\begin{aligned}
 \widehat V^+(x_k,u)&=G(1-\cos\widehat\alpha_N),\\
 \widehat\Theta(x_k,u)&=\max_{0\le j\le N}|\widehat\theta_j|,\\
 \mathcal J(x_k,u;w_k)&=\left(\frac{\Estar-\widehat V^+(x_k,u)}{\Estar}\right)^2\\
 &\quad+\lambda_W\left(\frac{R(u)}{W_{\max}}\right)^2\\
 &\quad+\lambda_\theta\left(\frac{\max\{\widehat\Theta(x_k,u)-\theta_{\max},0\}}{\theta_{\max}}\right)^2,\\
 u_k&\in\underset{u\in\mathcal U_{\mathrm{test}}}{\arg\min}\ \mathcal J(x_k,u;w_k).
\end{aligned}
\label{eq:score}
\end{equation}
In $\mathcal J(x_k,u;w_k)$, normalization by $\Estar$, $W_{\max}$, and $\theta_{\max}$ makes the three residuals dimensionless. The positive weights $\lambda_W,\lambda_\theta$ are tunable. Minimizing $\mathcal J(x_k,u;w_k)$ favors $\widehat V^+(x_k,u)=\Estar$, $R(u)=0$, and $\widehat\Theta(x_k,u)\le\theta_{\max}$. All $u\in\mathcal U_{\mathrm{test}}$ compete under $\mathcal J(x_k,u;w_k)$, so $R(u_k)$ need not vanish. After applying $u_k$, the controller observes $x_{k+1}$ and repeats Eqs.~\eqref{eq:workpolicy}--\eqref{eq:score}. Membership in $\mathcal U_{\mathrm{test}}$ enforces $|u_k|\le U$. A finite $\lambda_\theta$ does not guarantee $|\theta(t)|\le\theta_{\max}$, and stability of $\mathcal X_\star$ under minimization of $\mathcal J(x_k,u;w_k)$ is not established.

% Mechanical coupling determines the physical-energy region independently of actuation.
Mechanical coupling restricts $\mathcal E$ and hence $\mathcal S(e_k,w_k)$. Arm rotation moves the pendulum hinge, so the minimum $K_p$ depends on $K_a$ and $\alpha$. Subtracting $K_a$ from $K$ in Eq.~\eqref{eq:lagrangian} and writing $K_p$ as a nonnegative square in $\nu$ gives
\begin{equation}
\begin{aligned}
 K_a&=\tfrac12I_a\omega^2,\qquad \omega^2=2K_a/I_a,\\
 K_p&=K-K_a\\
    &=\tfrac12A\omega^2+C\omega\nu+\tfrac12J\nu^2
       -\tfrac12I_a\omega^2\\
    &=\tfrac12(A-I_a)\omega^2-\frac{C^2}{2J}\omega^2\\
    &\quad+\frac J2\left[\nu^2+\frac{2C}{J}\omega\nu
                    +\frac{C^2}{J^2}\omega^2\right]\\
    &=\frac J2\left(\nu+\frac CJ\omega\right)^2
       +\frac{A-I_a-C^2/J}{I_a}K_a,\\
 K_a&\ge0,\qquad 0\le V\le2G,\\
 K_p&\ge\frac{A-I_a-C^2/J}{I_a}\,K_a.
\end{aligned}
\label{eq:energy-feasibility}
\end{equation}
The term $\frac J2(\nu+C\omega/J)^2$ vanishes at $\nu=-C\omega/J$, giving the minimum $K_p$ for prescribed $K_a$ and $V$. Potential energy $V$ ranges from zero at the downward position to $2G$ at upright. Substituting $\cos\alpha=1-V/G$ into $A(\alpha)$ and $C(\alpha)$ expresses Eq.~\eqref{eq:energy-feasibility} entirely in energy coordinates,
\begin{equation}
\begin{aligned}
 C&=b(1-V/G),\\
 A&=I_a+m_pr^2+J\bigl[1-(1-V/G)^2\bigr].
\end{aligned}
\label{eq:energy-inertias}
\end{equation}
Thus the multiplier $(A-I_a-C^2/J)/I_a$ of $K_a$ depends on $V$ without requiring the sign of $\alpha$.

% Constructing real velocities proves realizability of every triple satisfying the mechanical bounds.
Every $e=(K_a,K_p,V)^{\mathsf T}$ satisfying Eq.~\eqref{eq:energy-feasibility} also corresponds to a physical state $x$ when velocities are unrestricted. Choose $\alpha$ with $\cos\alpha=1-V/G$, then choose $\omega$ supplying $K_a$. The velocity $\nu$ supplies $K_p$ above the minimum $(A-I_a-C^2/J)K_a/I_a$ through
\begin{equation}
\begin{aligned}
 \omega&=\pm\sqrt{2K_a/I_a},\\
 \nu&=-\frac CJ\omega\\
 &\quad\pm\sqrt{\frac2J\left[K_p-\frac{A-I_a-C^2/J}{I_a}K_a\right]}.
\end{aligned}
\label{eq:energy-realization}
\end{equation}
Equation~\eqref{eq:energy-feasibility} ensures real $\omega$ and $\nu$ in Eq.~\eqref{eq:energy-realization}, with either sign of $\omega$ permitted. Choose $\theta\in[-\theta_{\max},\theta_{\max}]$ independently because $H(x)$ does not depend on $\theta$. Hence Eqs.~\eqref{eq:energy-feasibility}--\eqref{eq:energy-inertias} characterize $\mathcal E$, but do not establish reachability from $x_k$ over $\Delta$. The inclusion $\mathcal C_{x_k}\subseteq\mathcal E$ holds through physical prediction, so the controller requires no additional filtering by Eq.~\eqref{eq:energy-feasibility}.


\section{Experiments and Results}
% R1: Single results paragraph; tables and atlas carry protocol and diagnostics.
We evaluated the selected controller in the nominal lossless simulator using four randomized initial-state classes (Table~\ref{tab:initial}). Parameter development used seeds 20261020 and 20261021. The reported setting was then evaluated on seeds 20261022 and 20261023, each supplying 512 starts per class. All 4,096 randomized episodes captured within 20 s (Fig.~\ref{fig:atlas}). The 1,024 downward starts remained inside the arm envelope and had a median capture time of 11.13 s. Six of the 1,024 moving starts crossed the arm boundary; the near-upright and tight-upright classes had no crossings. Eight fixed-probe episodes, counted separately, also captured without crossings. The replay traces show repeated energy exchange during swing-up and two observed crossing conditions. In one, the finite penalty permits a crossing despite bounded tested alternatives. In the other, no tested torque keeps the next action inside the envelope. These results support capture on the specified nominal distributions. Sustained holding, robustness to model error, and intervention savings require continued trajectories and a comparison study.












% \begin{table*}[t]
% \centering\small
% \caption{Physical inputs and selected controller settings. Physical values are configured model inputs; feedback and penalty values are design choices. The upright energy is derived from the model.}
% \label{tab:parameters}
% \begin{tabular}{lll}
% \toprule Quantity & Symbol & Value \\
% \midrule Arm mass; length & $m_a;\ r$ & $0.095$ kg; $0.085$ m \\
% Pendulum mass; length & $m_p;\ l$ & $0.024$ kg; $0.129$ m \\
% Gravity; viscous damping & $g$; both joints & $9.81$ m\,s$^{-2}$; zero \\
% Physical; usable torque limits & $u_{\mathrm{phys}};\ U$ & $0.0204$; $0.01836$ N\,m \\
% Desired arm envelope & $\theta_{\max}$ & $\pi$ rad \\
% Upright-rest energy & $\Estar=2G$ & $0.03037176$ J \\
% Held action; prediction step & $\Delta;\ h$ & $0.1$; $0.02$ s \\
% Work gain; request bound & $\kappa;\ W_{\max}$ & $0.004$; $0.04\Estar$ \\
% Work; arm penalty weights & $\lambda_W;\lambda_\theta$ & $0.02$; $10{,}000$ \\
% Torque grid; bisection steps & --- & $33$ points; $10$ steps \\
% \bottomrule
% \end{tabular}
% \end{table*}

\begin{table*}[t]
\centering\small
\caption{Validation initial states. Coordinates are sampled independently and uniformly within each stated interval; angles are in rad and velocities in rad\,s$^{-1}$. The moving pendulum angle receives an independent equiprobable sign. Angles are then wrapped to $[0,2\pi)$. Each randomized class contains 1,024 episodes across two seeds. The same four deterministic probes are repeated once per seed. Plant and controller predictions use 20 ms integration steps in these reported records.}
\label{tab:initial}
\begin{tabular}{lrrrr}
\toprule Class & $\theta$ & $\alpha$ & $\omega$ & $\nu$ \\
\midrule Downward & $[-0.20,0.20]$ & $[-0.20,0.20]$ & $[-0.5,0.5]$ & $[-0.5,0.5]$ \\
Moving & $[-0.50,0.50]$ & $\pm[0.40,2.60]$ & $[-2,2]$ & $[-6,6]$ \\
Near upright & $[-0.25,0.25]$ & $\pi+[-0.25,0.25]$ & $[-1,1]$ & $[-1,1]$ \\
Tight upright & $[-0.08,0.08]$ & $\pi+[-0.12,0.12]$ & $[-0.15,0.15]$ & $[-0.30,0.30]$ \\
\midrule
\multicolumn{5}{l}{Fixed probes: $(0,0,0,0)$, $(0,\pi,0,0)$, $(1.45,0.4,2,0)$, $(-1.45,-0.4,-2,0)$.}\\
\bottomrule
\end{tabular}
\end{table*}

\begin{thebibliography}{9}
\bibitem{heemels2012}
W. P. M. H. Heemels, K. H. Johansson, and P. Tabuada,
``An introduction to event-triggered and self-triggered control,''
in \emph{Proceedings of the 51st IEEE Conference on Decision and Control},
2012, pp. 3270--3285. \href{https://doi.org/10.1109/CDC.2012.6425820}{doi:10.1109/CDC.2012.6425820}.
\bibitem{astrom2000}
K. J. \AA str\"om and K. Furuta,
``Swinging up a pendulum by energy control,''
\emph{Automatica}, vol. 36, no. 2, 2000.
\href{https://doi.org/10.1016/S0005-1098(99)00140-5}{doi:10.1016/S0005-1098(99)00140-5}.
\end{thebibliography}
\end{document}
